\documentclass[11pt]{article}
\usepackage[margin=1in]{geometry}
\usepackage{amsmath,amssymb,amsthm,booktabs,array}
\usepackage{graphicx}
\usepackage[hidelinks]{hyperref}
\usepackage[T1]{fontenc}
\usepackage{microtype}

\newtheorem{theorem}{Theorem}
\newtheorem{lemma}[theorem]{Lemma}
\newcommand{\C}{\mathbb C}
\newcommand{\D}{\mathcal D}
\newcommand{\K}{\mathcal K}
\newcommand{\Z}{\mathcal Z}
\newcommand{\adj}{\operatorname{adj}}
\newcommand{\rank}{\operatorname{rank}}
\newcommand{\tr}{\operatorname{tr}}

\title{The degree-six polynomial in Helton--Klep--McCullough--Vol\v{c}i\v{c} is minimal\\
\large An exact computer-assisted full solution of the degree-five-versus-six question}
\author{Autonomous proof search, run \texttt{fa\_banach\_001}, lane 2}
\date{11 August 2026}

\begin{document}
\maketitle

\noindent\textbf{Status.} Claimed full solution, likely valid, pending expert review.  The finite algebraic classification in the proof is certified exactly over $\mathbb Q(\sqrt2)$ by the included verifier; no floating-point computation is used.

\section{Source question}

The source is J.~W.~Helton, I.~Klep, S.~McCullough and J.~Vol\v{c}i\v{c},
\emph{Noncommutative polynomials describing convex sets}, arXiv:1808.06669, Section~5.2, physical PDF page~24 \cite{HKMV}.

\begin{center}
\includegraphics[width=\textwidth]{figures/open_problem_crop.png}
\end{center}

Put $y=x^*$ and $H=x+y$.  The source defines
\begin{align*}
 f_1&=1-H-2H^2-2yx+H^3+2H^2yx,\\
 s&=1-H^2,
\end{align*}
and
\[
L=\begin{pmatrix}
1-\frac12H&-\sqrt2H&\frac12H&y\\
-\sqrt2H&1&0&0\\
\frac12H&0&1-\frac12H&-y\\
x&0&-x&1
\end{pmatrix}.
\]
It proves that $f=f_1s$ is Hermitian, $\D_f=\D_L$, $\deg f=6$, and every scalar Hermitian defining polynomial for this domain has degree at least five.  The exact source sentence is:

\begin{quote}
``Note that $\deg f=6$, but we do not know whether $f$ is a minimal degree defining polynomial.''
\end{quote}

Thus the missing issue is the existence or nonexistence of a degree-five scalar Hermitian polynomial defining $\D_L$.

\section{Definitions and input results}

We work in $\C\langle x,y\rangle$ with involution $x^*=y$, $y^*=x$.  For a scalar polynomial $q$ with $q(0)=1$, $\K_q$ is the closure of the connected component of the origin in the set of matrix tuples where $q$ is invertible.  When $q=q^*$, the same set is denoted $\D_q$.  The free singularity locus is $\Z_q$.

We use the following results exactly as they are used in the source paper.
\begin{enumerate}
\item Lemma~5.1 of \cite{HKMV}, together with the Polynomial Singularit\"atstellensatz, says that any Hermitian $q$ with $\D_q=\D_L$ has degree at least five; moreover, a degree-five such $q$ has an atomic factor of degree four in the stable-association class of $f_1$ \cite[Theorem~4.3]{HKV}.
\item Proposition~2.7 and Remark~2.6 of \cite{HKMV} give the minimal Fornasini--Marchesini realization, the determinant identity, the inversion formula, and the joint-nilpotence criterion for a realization to be a polynomial.
\item The proof of Proposition~3.3 of \cite{HKMV} applies verbatim to an arbitrary normalized atom $p$ with $\Z_p=\Z_L$: a minimal realization of $p^{-1}$ can be written over the displayed pencil $L$.  Indeed, its irreducible realization pencil has the same free locus as $L$ and hence is similar to $L$ by \cite[Theorem~3.11]{KV}.
\end{enumerate}

\section{Proof intuition}

The source lower bound leaves only one possible way a degree-five defining polynomial could exist: it would have to be a degree-four atom with the same free singularity locus as $f_1$, multiplied by an affine-linear factor.  Stable association means that the degree-four atom need not literally be $f_1$, so checking only $f_1t$ is insufficient.

Minimal realization theory makes all stable associates finite in this example.  A degree-four atom $p$ with $\Z_p=\Z_L$ gives a rank-one matrix $M=v\alpha$ satisfying the scalar adjugate identity
\[
\tr(M\adj L(z,\bar z))=1.
\]
Coefficient comparison leaves a seven-dimensional affine space, and the rank-one minors cut it down to exactly four points.  These yield four explicit degree-four atoms (one is $f_1$).  For each atom, six coefficients of $pt-(pt)^*$ form a nonsingular linear system in the six real and imaginary parts of an affine $t$.  Hence $t=0$, contradicting a degree-five factorization.

\section{Main theorem}

\begin{theorem}\label{thm:main}
Let $f_1,s,L$ be the polynomials and pencil above and set $f=f_1s$.  If $q\in\C\langle x,y\rangle$ is Hermitian, $q(0)=1$, and $\D_q=\D_L$, then $\deg q\ge 6$.  Consequently $f$ is a minimal-degree scalar Hermitian defining polynomial for $\D_L$, and the minimum degree is exactly six.
\end{theorem}

\begin{proof}
The source paper proves $f=f^*$, $\D_f=\D_L$, and $\deg f=6$, as well as the lower bound $\deg q\ge5$ \cite[Lemma~5.1 and Section~5.2]{HKMV}.  It remains to exclude degree five.

\paragraph{1. Reduction to a linear symmetrizer.}
Suppose $q=q^*$, $\D_q=\D_L$, and $\deg q=5$.  Normalize every factor to have constant term one.  The argument of Lemma~5.1 in \cite{HKMV}, using \cite[Theorem~4.3(3)]{HKV}, gives an atomic factor $p$ of degree four with
\[
\Z_p=\Z_{f_1}=\Z_L.
\]
Degree is additive in the free algebra, so all remaining atomic factors have total degree one.  If $p$ occurs on the left, $q=pt$ for a nonzero affine-linear $t$.  If it occurs on the right, take adjoints and use $q=q^*$ to obtain the same form $q=p^*t^*$; here $p^*$ has the same free singularity locus.  Thus it suffices to prove:
\begin{equation}\label{eq:goal}
\text{no normalized degree-four atom $p$ with $\Z_p=\Z_L$ has $pt=(pt)^*$ for nonzero affine $t$.}
\end{equation}

\paragraph{2. Realization of every possible $p$.}
Write
\[
L=I-A_xx-A_yy,
\quad
A_x=\begin{pmatrix}
\frac12&\sqrt2&-\frac12&0\\
\sqrt2&0&0&0\\
-\frac12&0&\frac12&0\\
-1&0&1&0
\end{pmatrix},
\quad A_y=A_x^{\mathsf T}.
\]
By the realization results recalled above, every $p$ in \eqref{eq:goal} has a minimal realization
\[
p^{-1}=1+\alpha L^{-1}(b_xx+b_yy),
\]
where $\alpha$ is a row vector.  The inversion formula says that the realization coefficients for $p$ are
\[
N_x=A_x-b_x\alpha,
\qquad
N_y=A_y-b_y\alpha.
\]
Since $p$ is a polynomial, $N_x,N_y$ are jointly nilpotent \cite[Remark~2.6(5)--(6)]{HKMV}.  Choose a nonzero common kernel vector $v$.  If $\alpha v=0$, then $A_xv=A_yv=0$, contradicting minimality by Remark~2.6(1) of \cite{HKMV}.  Scale $v$ so that $\alpha v=1$.  Then
\[
b_x=A_xv,\qquad b_y=A_yv,\qquad
p^{-1}=\alpha L^{-1}v.
\]
Proposition~2.7(1) of \cite{HKMV}, evaluated at commuting scalar variables $z,\bar z$, gives $p(z,\bar z)=\det L(z,\bar z)$.  Hence
\begin{equation}\label{eq:adj}
\alpha\adj L(z,\bar z)v=1.
\end{equation}
Set $M=v\alpha$.  Then $M$ has rank one and \eqref{eq:adj} is
\begin{equation}\label{eq:trace}
\tr\!\left(M\adj L(z,\bar z)\right)=1.
\end{equation}

\paragraph{3. Exact rank-one classification.}
Equating the coefficients in \eqref{eq:trace} gives every solution in the form
\[
M(a,b,c,d,e,r,g)=
\begin{pmatrix}
1-b+r/\sqrt2&e/\sqrt2+g/(2\sqrt2)-a&g-b-c-1&e\\
a&b&-e/\sqrt2-g/(2\sqrt2)-d&r\\
c&d&-r/\sqrt2-g&g/2\\
e&r&g/2&g
\end{pmatrix}.
\]
The lexicographic Gr\"obner basis of its $2\times2$ minors contains
\begin{gather*}
e, r, g, d^2, cd, c^2(c+1), c(b+c+1), d(b+1),\\
a c+2d,\quad b^2+b-2c^2-2c,\quad ab+a-4d,\quad
a^2+2b-2c^2-2c.
\end{gather*}
The complete basis is printed by the verifier.  These equations force $d=e=r=g=0$.  If $c=-1$, they give $(a,b)=(0,0)$; if $c=0$, they give $(a,b)=(0,0)$ or $(\pm\sqrt2,-1)$.  Substitution verifies that these four points satisfy all minors.  Therefore the only rank-one matrices are
\begin{align*}
M_1&=\begin{pmatrix}1&0&0&0\\0&0&0&0\\-1&0&0&0\\0&0&0&0\end{pmatrix},&
M_2&=\begin{pmatrix}1&0&-1&0\\0&0&0&0\\0&0&0&0\\0&0&0&0\end{pmatrix},\\[1ex]
M_3&=\begin{pmatrix}2&\sqrt2&0&0\\-\sqrt2&-1&0&0\\0&0&0&0\\0&0&0&0\end{pmatrix},&
M_4&=\begin{pmatrix}2&-\sqrt2&0&0\\\sqrt2&-1&0&0\\0&0&0&0\\0&0&0&0\end{pmatrix}.
\end{align*}
Choose factorizations $M_i=v_i\alpha_i$ as follows:
\[
\begin{array}{c|c|c}
i&v_i&\alpha_i\\ \hline
1&(1,0,-1,0)^{\mathsf T}&(1,0,0,0)\\
2&(1,0,0,0)^{\mathsf T}&(1,0,-1,0)\\
3&(\sqrt2,-1,0,0)^{\mathsf T}&(\sqrt2,1,0,0)\\
4&(\sqrt2,1,0,0)^{\mathsf T}&(\sqrt2,-1,0,0).
\end{array}
\]
For $j\in\{x,y\}$ put
\[
N_{i,j}=A_j(I-M_i),\qquad b_{i,j}=A_jv_i,\qquad
N_i=N_{i,x}x+N_{i,y}y.
\]
Direct exact multiplication gives $N_{i,w}=0$ for every word $w$ of length four.  Thus the corresponding polynomial is
\begin{equation}\label{eq:pi}
p_i=1-\alpha_i(I+N_i+N_i^2+N_i^3)(b_{i,x}x+b_{i,y}y).
\end{equation}
Each $p_i$ has degree four, and $p_2=f_1$.  Equations \eqref{eq:trace}--\eqref{eq:pi} show that every possible $p$ is one of these four.

\paragraph{4. No affine symmetrizer.}
Let
\[
t=u_0+u_xx+u_yy,\qquad
u_0=u_{0r}+iu_{0i},\quad u_x=u_{xr}+iu_{xi},\quad u_y=u_{yr}+iu_{yi}.
\]
For each $i$, equate coefficients in $p_it-(p_it)^*=0$.  The following six coefficients already have full rank.  A word such as $xxxyx$ denotes the corresponding word in the free letters $x,y$.

\begin{center}\small
\renewcommand{\arraystretch}{1.15}
\begin{tabular}{c|l}
\toprule
$i$&Six coefficient equations\\
\midrule
1&$-2u_{yr}=0$, $2u_{yi}=0$, $2u_{xr}=0$, $2u_{xi}=0$, $-2u_{0r}+u_{xr}-u_{yr}=0$, $2u_{0i}+u_{xi}+u_{yi}=0$\\
2&$2u_{xr}=0$, $2u_{xi}=0$, $2u_{yr}=0$, $2u_{yi}=0$, $2u_{0r}+u_{xr}-u_{yr}=0$, $2u_{0i}+u_{xi}+u_{yi}=0$\\
3&$-2u_{yr}=0$, $2u_{yi}=0$, $2u_{xr}=0$, $2u_{xi}=0$, $4u_{0i}-u_{xi}-u_{yi}=0$, $-4u_{0r}-2u_{xr}+4u_{yr}=0$\\
4&$-2u_{yr}=0$, $2u_{yi}=0$, $2u_{xr}=0$, $2u_{xi}=0$, $4u_{0i}+3u_{xi}+3u_{yi}=0$, $4u_{0r}-2u_{xr}+4u_{yr}=0$\\
\bottomrule
\end{tabular}
\end{center}

For transparency, the selected words and real/imaginary parts are printed by the verifier.  The determinants of these four $6\times6$ systems are respectively
\[
64,\qquad64,\qquad-256,\qquad256.
\]
Therefore $u_0=u_x=u_y=0$ for every $i$.  This proves \eqref{eq:goal}, contradicting the nonzero affine factor in a hypothetical degree-five $q$.

Hence every Hermitian defining polynomial has degree at least six.  Since the source polynomial $f=f_1s$ has degree six and defines $\D_L$, it is minimal.  \qedhere
\end{proof}

\section{Exact verification}

The included script \texttt{code/verifier.py} performs the following exact checks over $\mathbb Q(\sqrt2)$:
\begin{enumerate}
\item $f_1s$ is Hermitian and has degree six;
\item the affine parametrization solves \eqref{eq:trace};
\item all $36$ rank-one minors have exactly the four solutions above, with the displayed Gr\"obner basis;
\item all words of length four in every pair $N_{i,x},N_{i,y}$ vanish;
\item every $p_i$ has degree four and $p_2=f_1$;
\item the coefficient matrix for $p_it-(p_it)^*$ has rank six, with the four nonzero determinants displayed above.
\end{enumerate}
Run it from the repository root with
\begin{verbatim}
conda run --no-capture-output -n sandbox python \
  runs/fa_banach_001/solutions/full/\
  1808.06669_degree_six_minimal_defining_polynomial/code/verifier.py
\end{verbatim}
The recorded output ends with
\begin{verbatim}
PASS: exactly four realization candidates; no affine symmetrizer exists.
\end{verbatim}
The computation is part of the proof only as an exact finite polynomial calculation; it is not numerical evidence.

\section{Limitations, novelty check, and review recommendation}

This result concerns the single explicit scalar example in Section~5.2 of \cite{HKMV}; it does not classify minimal defining degrees for arbitrary free spectrahedra.  The proof relies on the cited realization and factorization theorems and on an exact finite Gr\"obner-basis/coefficient-rank certificate.

On 11 August 2026, a bounded novelty search checked the run's registry, solution, attempt, and proof-gap indexes and searched arXiv by the exact source sentence, the phrases ``minimal degree defining polynomial'' and ``free spectrahedron,'' the title and arXiv id 1808.06669, and combinations of the authors' names with the degree-five-versus-six terminology.  Only the source paper and background factorization papers were found; no later paper explicitly resolving this example was located.  This is evidence of novelty, not a guarantee of priority.

Human review should focus on two points: (i) the passage from a degree-five Hermitian defining polynomial to a degree-four atomic factor realized over $L$, and (ii) the completeness of the four-point rank-one classification.  Both are isolated explicitly above, and the second has a rerunnable exact certificate.

\begin{thebibliography}{9}
\bibitem{HKMV}
J.~W. Helton, I. Klep, S. McCullough and J. Vol\v{c}i\v{c},
\emph{Noncommutative polynomials describing convex sets},
arXiv:1808.06669; J. Funct. Anal. 279 (2020), 108324.

\bibitem{HKV}
J.~W. Helton, I. Klep and J. Vol\v{c}i\v{c},
\emph{Geometry of free loci and factorization of noncommutative polynomials},
Adv. Math. 331 (2018), 589--626; arXiv:1708.05378.

\bibitem{KV}
I. Klep and J. Vol\v{c}i\v{c},
\emph{Free loci of matrix pencils and domains of noncommutative rational functions},
Comment. Math. Helv. 92 (2017), 105--130.
\end{thebibliography}

\end{document}
